\section{Fixed bounded formulas and Boolean reflection}
\label{sec:boolean}

The induction scheme of \iPRA{} applies only to quantifier-free formulas.  This is not an obstacle to finite mathematics, provided bounded logical content is compiled to primitive-recursive bits before induction.  The word \emph{fixed} is important: compilation is structural recursion in Lean on a theorem-specific syntax tree, not evaluation of an arbitrary object-language program.

\subsection{A fixed bounded grammar}

For a fixed number of free variables, consider formulas generated by
\[
  A ::= s=t\mid s<t\mid A\wedge A\mid A\vee A\mid A\to A
       \mid \exists i<N\,A(i)\mid \forall i<N\,A(i),
\]
where all terms are fixed primitive-recursive terms.  Negation is \(A\to\bot\).  Bounds may depend on the outer free variables but not on the variable they bind.

\begin{theorem}[Boolean reflection]
\label{thm:boolean-reflection}
For every fixed bounded formula \(A(\vec x)\), one can effectively construct a primitive-recursive term \(\chi_A(\vec x)\) and object derivations, under every context \(\Gamma\), of
\begin{align}
 &\chi_A=0\;\vee\;\chi_A=1, \label{eq:reflect-bit}\\
 &\chi_A=1\;\leftrightarrow\; A, \label{eq:reflect-one}\\
 &\chi_A=0\;\leftrightarrow\; \neg A. \label{eq:reflect-zero}
\end{align}
The construction is uniform at the metalevel in the syntax tree of \(A\).  It introduces no new nonlogical axiom.
\end{theorem}

\begin{proof}
We proceed by structural recursion on \(A\).

\paragraph{Equality.}
Take \(\chi_{s=t}=\mathsf{eqb}(s,t)\).  The checked equality-bit development proves \(\mathsf{eqb}(s,t)=1\) from \(s=t\), and \(\mathsf{eqb}(s,t)=0\) from \(s\neq t\).  The checked totality of equality gives \(s=t\vee s\neq t\), so the converse directions and bit normalization follow by disjunction elimination.

\paragraph{Order.}
Use the characteristic bit for the repository's chosen order definition.  Order is expressed through monus and equality, so its normalization and reflection reduce to the preceding case plus the checked order theorems.

\paragraph{Conjunction.}
Let \(b=\chi_B\), \(c=\chi_C\), and set \(\chi_{B\wedge C}=bc\).  Four finite cases from \eqref{eq:reflect-bit} for \(b,c\) establish normalization.  If \(bc=1\), neither factor can be zero; the Boolean cases force \(b=c=1\), and the induction hypotheses give \(B\wedge C\).  Conversely, \(B\wedge C\) yields \(b=c=1\), hence \(bc=1\).  The zero case is similar and intuitionistically valid because both component formulas are decidable.

\paragraph{Disjunction.}
Set \(\chi_{B\vee C}=\mathsf{sg}(b+c)\), where \(\mathsf{sg}(0)=0\) and \(\mathsf{sg}(S n)=1\).  If one input bit is one, the sum is a successor and the output bit is one.  If the output is one, the sum is nonzero; finite Boolean case analysis excludes \(b=c=0\), so one component formula holds.  The zero direction follows from \(b+c=0\), which implies both summands are zero.

\paragraph{Implication.}
Define
\[
  \chi_{B\to C}=\mathsf{or}(\mathsf{not}(b),c).
\]
Classically one writes \(B\to C\equiv\neg B\vee C\).  Here that equivalence is justified intuitionistically because \(B\) is decidable by the induction hypothesis: split on \(B\vee\neg B\).  In the first branch an implication gives \(C\); in the second branch \(\neg B\) is already available.  The reverse direction is intuitionistically immediate.  The Boolean truth table then gives \eqref{eq:reflect-one}--\eqref{eq:reflect-zero}.

\paragraph{Bounded existential.}
Assume \(B(i,\vec x)\) has reflected bit \(b(i,\vec x)\).  Define a finite sum
\[
  S(N,\vec x)=\sum_{i<N}b(i,\vec x)
\]
and put \(\chi_{\exists i<N B}=\mathsf{sg}(S)\).  The defining equations of the checked finite-sum symbol are available at arbitrary terms.  Quantifier-free induction on \(N\) proves
\begin{align*}
  S(N,\vec x)=0 &\;\to\; \forall i<N\;b(i,\vec x)=0,\\
  \bigl(\exists i<N\;b(i,\vec x)=1\bigr) &\;\to\; S(N,\vec x)\neq 0.
\end{align*}
For witness extraction from \(S\neq 0\), one may either use the same induction or, more modularly, invoke bounded least search on the zero-body \(1\dot{-} b(i,\vec x)\).  The search theorem returns the least index with \(b=1\).  Reflection for \(B\) converts between the bit and formula witnesses.

\paragraph{Bounded universal.}
There are two convenient definitions.  One may use the product
\[
  P(N,\vec x)=\prod_{i<N}b(i,\vec x)
\]
after proving that all factors are bits, or define the universal bit as the negation of the existential failure bit:
\[
  \chi_{\forall i<N B}
  =\mathsf{not}\bigl(\chi_{\exists i<N\neg B}\bigr).
\]
The second choice reuses the existential development.  Because every instance \(B(i,\vec x)\) is decidable, intuitionistic bounded De Morgan reasoning proves
\[
 \neg\exists i<N\neg B(i)\;\leftrightarrow\;\forall i<N B(i).
\]
The proof is a quantifier-free induction on \(N\), or a bounded least-counterexample argument.

Each constructor preserves substitution because the generated terms and formulas are built from syntax whose substitution laws are already proved.  Thus the result is usable under arbitrary local hypotheses and at arbitrary terms.
\end{proof}

\subsection{Least witnesses and counterexamples}

Reflection is more useful when it comes with extraction.

\begin{corollary}[Least bounded witness]
\label{cor:least-bounded-witness}
For every fixed reflected body \(B(i,\vec x)\), there is a primitive-recursive term \(\mu_B(N,\vec x)\) such that \iPRA{} derives:
\begin{enumerate}[label=(\roman*),leftmargin=2.3em]
\item \(\mu_B\leq N\);
\item if \(\exists i<N\,B(i,\vec x)\), then \(\mu_B<N\), \(B(\mu_B,\vec x)\), and \(\mu_B\leq i\) for every supplied witness \(i<N\);
\item if \(j<\mu_B\), then \(\neg B(j,\vec x)\).
\end{enumerate}
\end{corollary}

\begin{proof}
Apply the checked \texttt{Search.mu} to the term
\[
  f(i,\vec x)=1\dot{-}\chi_B(i,\vec x).
\]
For normalized bits, \(f=0\) is equivalent to \(\chi_B=1\), hence to \(B\).  The statements are direct instances of \texttt{search\_spec}, \texttt{search\_found}, and \texttt{search\_minimal}, followed by the reflection equivalence.
\end{proof}

The dual least-counterexample term is obtained by searching for \(\chi_B=0\).  This is the basis of finite failure certificates and universal strategy checking.

\subsection{Fixed alternation and finite strategies}

A fixed bounded quantifier prefix
\[
  Q_0x_0<B_0\;Q_1x_1<B_1\cdots Q_{m-1}x_{m-1}<B_{m-1}\;A(\vec x)
\]
can be evaluated backwards as a finite game.  Existential levels take a finite disjunction; universal levels take a finite conjunction.  A success certificate stores one move at existential levels and all required child certificates at universal levels.  A failure certificate reverses the roles.  Since \(m\) and the source formula are fixed, the resulting verifier is primitive recursive and its soundness is proved by induction on the fixed syntax tree plus quantifier-free induction over variable bounds.

This recovers the useful part of finite classical reasoning: every bounded position has a decidable winning bit and one of two finite strategy certificates.  It does \emph{not} recover an unbounded choice function or a witness for an arbitrary classically proved existential statement.

\subsection{Lean implementation route}

A practical Lean implementation should define a small metalevel datatype \texttt{BFormula n} with constructors for atoms, connectives, and bounded quantifiers.  Compilation returns a structure containing a \texttt{PTm n}, an \texttt{Fm PTm n}, and the five derivations in \cref{thm:boolean-reflection}.  The constructors should call existing arithmetic lemmas rather than reprove Boolean algebra.  In particular:

\begin{itemize}[leftmargin=2em]
\item equality uses \texttt{eqb\_of\_eq}, \texttt{eqb\_of\_ne}, and \texttt{eq\_or\_ne};
\item bounded witnesses use \texttt{Search.search\_found};
\item folds use \texttt{sumF}/\texttt{prodF} and their equations;
\item the only induction is \texttt{LeanProofs.IPRA.ind} on a quantifier-free accumulator invariant.
\end{itemize}

The generated \texttt{PR} symbol is specific to the fixed \texttt{BFormula}.  There is no object-level syntax code or evaluator unless a particular application independently requires one.
